\section{Introduction}

The free uniform spanning forest is a central probability measure on infinite
graphs; see, for example, \cite{BLPS2001,LyonsPeres2016}.  A recent OpenAI
preprint proves that the free uniform spanning forest on every infinite,
connected, locally finite, simple, unweighted graph is a factor of independent
vertex labels \cite{OpenAI2026FUSF}.  A principal ingredient is a Brownian
sampler for binary processes satisfying a uniform response bound.  Its output
is recovered from the asymptotic slopes of a deterministically reconstructed
observation field.

The sampler proof contains a stronger path-law statement than is needed merely
to identify the final output.  If $W$ denotes the independent Brownian input,
$Z(W)$ the reconstructed path field, and $\Psi(W)$ the binary output, then
\begin{equation}\label{eq:path-law-intro}
  \bigl(Z(W),\Psi(W)\bigr)
  \ \stackrel{\mathrm d}{=}\
  \bigl((tX_i+B_i(t))_{i,t},X\bigr),
\end{equation}
where $X$ has the target law and the Brownian motions $B_i$ are independent of
one another and of $X$.  We take \eqref{eq:path-law-intro} as the sole input
from the sampler construction and determine exactly how quickly the natural
midpoint readout stabilizes.

The distinction between a snapshot error and a future error is important.  At
a fixed time $t$, one coordinate is decoded incorrectly with probability
$\Phi(-\sqrt t/2)$.  The probability that it makes an error at some time after
$t$ is twice as large:
\begin{equation}\label{eq:factor-two-intro}
  \P(T_i>t)=2\Phi\!\left(-\frac{\sqrt t}{2}\right).
\end{equation}
The second half of this probability comes from paths that are currently on the
correct side of the threshold but later return to it.  Thus $T_i$ is a
last-passage variable rather than a stopping time.

Our main theorem identifies the entire marked law:
\begin{equation}\label{eq:product-marking-intro}
  \Law\bigl(X,(T_i)_{i\in I}\bigr)
  =\Law(X)\otimes
  \bigotimes_{i\in I}\GammaLaw\!\left(\frac12,8\right).
\end{equation}
Arbitrary dependence in the target process disappears from the stabilization
marks.  The reason is the midpoint threshold: after a sign change determined
by $X_i$, either target bit produces the same Brownian motion with drift
$-1/2$.

The proof of the stabilization laws is self-contained after
\eqref{eq:path-law-intro}.  We record a Brownian last-passage lemma, derive the
product marking, and then develop exact finite-size, extreme-value, and
heterogeneous-clock consequences.  For finite target laws we additionally use
the posterior-mean interpretation of the sampler drift and the vector
I--MMSE identity to record an exact entropy-dissipation formula.  Standard
background on Brownian hitting and last-passage identities may be found in
\cite{KaratzasShreve1991,RevuzYor1999,BorodinSalminen2002}; background on
extreme-value limits may be found in \cite{Resnick1987}.

The stabilization parameter studied here is Brownian observation time after
the path field $Z(W)$ has been reconstructed.  It is not a coding radius in
the underlying IID vertex labels, and none of the results below asserts that
the factor map is finitary.  The exact law instead quantifies the online
readout of the reconstructed observation field.

\subsection*{Provenance}
The deterministic inversion, innovation representation, posterior-drift
identification, and path-law identity \eqref{eq:path-law-intro} are due to the
OpenAI preprint \cite{OpenAI2026FUSF}, cited at the fixed source commit recorded
in the bibliography.  The stabilization theorems and their consequences below
are derived from the path-law identity; the entropy formula also uses the
posterior-drift identification.
